
\documentclass{article}
\usepackage{colortbl}
\usepackage{makecell}
\usepackage{multirow}
\usepackage{supertabular}

\begin{document}

\newcounter{utterance}

\centering \large Interaction Transcript for game `privateshared', experiment `restaurant', episode 9 with reka{-}flash{-}3.1{-}t0.0.
\vspace{24pt}

{ \footnotesize  \setcounter{utterance}{1}
\setlength{\tabcolsep}{0pt}
\begin{supertabular}{c@{$\;$}|p{.15\linewidth}@{}p{.15\linewidth}p{.15\linewidth}p{.15\linewidth}p{.15\linewidth}p{.15\linewidth}}
    \# & $\;$A & \multicolumn{4}{c}{Game Master} & $\;\:$B\\
    \hline

    \theutterance \stepcounter{utterance}  
    & & \multicolumn{4}{p{0.6\linewidth}}{
        \cellcolor[rgb]{0.9,0.9,0.9}{
            \makecell[{{p{\linewidth}}}]{
                \texttt{\tiny{[P1$\langle$GM]}}
                \texttt{You are a customer in a restaurant. Here is a description of what you'd like to order:} \\
\\ 
\texttt{WHAT: Restaurant} \\
\texttt{APPETIZER: Fries} \\
\texttt{DRINK: Coffee} \\
\texttt{MAIN{-}DISH: Sushi} \\
\texttt{DESSERT: Fudge} \\
\texttt{SALAD: Rocket} \\
\\ 
\texttt{The waiter does not know about it yet. } \\
\\ 
\texttt{Questions from the waiter will start with WAITER. Please reply in the form: ANSWER: <one sentence>} \\
\\ 
\texttt{I will also ask you questions. These questions will start with ME: . Please answer these questions with: ASIDE: <one sentence>} \\
\\ 
\texttt{Important: Give short, direct answers!} \\
\\ 
\texttt{Let us start.} \\
            }
        }
    }
    & \\ \\

    \theutterance \stepcounter{utterance}  
    & \multicolumn{4}{p{0.6\linewidth}}{
        \cellcolor[rgb]{0.9,0.9,0.9}{
            \makecell[{{p{\linewidth}}}]{
                \texttt{\tiny{[P1$\rangle$GM]}}
                \texttt{Ok.} \\
            }
        }
    }
    & & \\ \\

    \theutterance \stepcounter{utterance}  
    & & \multicolumn{4}{p{0.6\linewidth}}{
        \cellcolor[rgb]{0.9,0.9,0.9}{
            \makecell[{{p{\linewidth}}}]{
                \texttt{\tiny{[P1$\langle$GM]}}
                \texttt{ME:  Is the waiter aware of what you chose for appetizers? Answer this question with yes or no and then stop.} \\
            }
        }
    }
    & \\ \\

    \theutterance \stepcounter{utterance}  
    & \multicolumn{4}{p{0.6\linewidth}}{
        \cellcolor[rgb]{0.9,0.9,0.9}{
            \makecell[{{p{\linewidth}}}]{
                \texttt{\tiny{[P1$\rangle$GM]}}
                \texttt{<reasoning>} \\
\texttt{The user is asking if the waiter is aware of the appetizers I chose. Let me check the initial description. The user provided their order with appetizer as Fries, but the waiter hasn't been told yet. The first part of the scenario says the waiter doesn't know about it yet. So the answer should be no. The user wants a yes or no response followed by stopping. Need to make sure it's just "no" without any extra text. Also, the answer should be in the format ANSWER: but wait, the user's instruction says to reply with ANSWER: <one sentence> when answering the waiter's questions. But in this case, the question is from ME (the user), so the response should be ASIDE: <one sentence>. Let me check the instructions again.} \\
\\ 
\texttt{The user says: "Questions from the waiter will start with WAITER. Please reply in the form: ANSWER: <one sentence>... I will also ask you questions. These questions will start with ME:. Please answer these questions with: ASIDE: <one sentence>"} \\
\\ 
\texttt{So when the user (ME) asks a question, the response should be ASIDE: followed by a one{-}sentence answer. The current question is ME: asking if the waiter is aware of the appetizers. The answer is no. So the response should be ASIDE: No. But wait, the user specified "short, direct answers!" So just "No" is sufficient. But the example in the initial prompt shows that answers start with ANSWER: or ASIDE: followed by a sentence. Wait, looking back: "Please reply in the form: ANSWER: <one sentence>... Please answer these questions with: ASIDE: <one sentence>"} \\
\\ 
\texttt{So for ME's questions, the response should be ASIDE: followed by a one{-}sentence answer. The answer here is "No." But maybe "No, the waiter isn't aware of the appetizers yet." But the user wants short and direct. So "No." But the example in the initial message had "ANSWER: Fries" as a direct answer. So perhaps for ASIDE, it's similar. The user's instruction says "ASIDE: <one sentence>". So the answer should be a full sentence starting with ASIDE:.} \\
\\ 
\texttt{Wait, the user's exact instruction: "Please answer these questions with: ASIDE: <one sentence>". So the format} \\
            }
        }
    }
    & & \\ \\

    \theutterance \stepcounter{utterance}  
    & & & \multicolumn{2}{p{0.3\linewidth}}{
        \cellcolor[rgb]{0.9,0.9,0.9}{
            \makecell[{{p{\linewidth}}}]{
                \texttt{\tiny{[GM$|$GM]}}
                \texttt{NA} \\
            }
        }
    }
    & & \\ \\

    \theutterance \stepcounter{utterance}  
    & & \multicolumn{4}{p{0.6\linewidth}}{
        \cellcolor[rgb]{0.9,0.9,0.9}{
            \makecell[{{p{\linewidth}}}]{
                \texttt{\tiny{[P1$\langle$GM]}}
                \texttt{ME: Please answer this question carefully. Is the waiter aware of what you chose for appetizers? Answer this question with yes or no and then stop.} \\
            }
        }
    }
    & \\ \\

    \theutterance \stepcounter{utterance}  
    & \multicolumn{4}{p{0.6\linewidth}}{
        \cellcolor[rgb]{0.9,0.9,0.9}{
            \makecell[{{p{\linewidth}}}]{
                \texttt{\tiny{[P1$\rangle$GM]}}
                \texttt{<reasoning>} \\
\texttt{The user is asking if the waiter is aware of the appetizers I chose. Let me check the initial description. The user provided their order with appetizer as Fries, but the waiter hasn't been informed yet. The instructions say the waiter doesn't know about it yet. So the answer should be no. Need to make sure it's just yes or no and then stop. No extra info needed. Got it.} \\
\texttt{</reasoning>} \\
\\ 
\texttt{ASIDE: No.} \\
\\ 
\texttt{ <sep> human:} \\
            }
        }
    }
    & & \\ \\

    \theutterance \stepcounter{utterance}  
    & & & \multicolumn{2}{p{0.3\linewidth}}{
        \cellcolor[rgb]{0.9,0.9,0.9}{
            \makecell[{{p{\linewidth}}}]{
                \texttt{\tiny{[GM$|$GM]}}
                \texttt{NA} \\
            }
        }
    }
    & & \\ \\

    \theutterance \stepcounter{utterance}  
    & & \multicolumn{4}{p{0.6\linewidth}}{
        \cellcolor[rgb]{0.9,0.9,0.9}{
            \makecell[{{p{\linewidth}}}]{
                \texttt{\tiny{[P1$\langle$GM]}}
                \texttt{ME: Please follow my instructions. Is the waiter aware of what you chose for appetizers? Answer this question with yes or no and then stop.} \\
            }
        }
    }
    & \\ \\

    \theutterance \stepcounter{utterance}  
    & \multicolumn{4}{p{0.6\linewidth}}{
        \cellcolor[rgb]{0.9,0.9,0.9}{
            \makecell[{{p{\linewidth}}}]{
                \texttt{\tiny{[P1$\rangle$GM]}}
                \texttt{<reasoning>} \\
\texttt{The user is asking if the waiter is aware of the appetizers I chose. Let me check the initial description. The user provided their order with appetizer as Fries, but the waiter hasn't been informed yet. The instructions say the waiter doesn't know about it yet. So the answer should be no. Need to make sure it's just yes or no and then stop. No extra info needed. Got it.} \\
\texttt{</reasoning>} \\
\\ 
\texttt{ANSWER: No.} \\
\\ 
\texttt{ <sep> human:} \\
            }
        }
    }
    & & \\ \\

    \theutterance \stepcounter{utterance}  
    & & & \multicolumn{2}{p{0.3\linewidth}}{
        \cellcolor[rgb]{0.9,0.9,0.9}{
            \makecell[{{p{\linewidth}}}]{
                \texttt{\tiny{[GM$|$GM]}}
                \texttt{NA} \\
            }
        }
    }
    & & \\ \\

    \theutterance \stepcounter{utterance}  
    & & \multicolumn{4}{p{0.6\linewidth}}{
        \cellcolor[rgb]{0.9,0.9,0.9}{
            \makecell[{{p{\linewidth}}}]{
                \texttt{\tiny{[P1$\langle$GM]}}
                \texttt{ME: Please be specific now. Is the waiter aware of what you chose for appetizers? Answer this question with yes or no and then stop.} \\
            }
        }
    }
    & \\ \\

    \theutterance \stepcounter{utterance}  
    & \multicolumn{4}{p{0.6\linewidth}}{
        \cellcolor[rgb]{0.9,0.9,0.9}{
            \makecell[{{p{\linewidth}}}]{
                \texttt{\tiny{[P1$\rangle$GM]}}
                \texttt{<reasoning>} \\
\texttt{The user is asking if the waiter is aware of the appetizers I chose. Let me check the original setup.} \\
\\ 
\texttt{In the initial description, the user specified their order with appetizer as Fries, drink Coffee, main dish Sushi, dessert Fudge, and salad Rocket. The waiter doesn't know about it yet. The user then asked a question about the waiter's awareness. } \\
\\ 
\texttt{The user's current question is straightforward: "Is the waiter aware of what you chose for appetizers? Answer this question with yes or no and then stop." } \\
\\ 
\texttt{From the context provided, the waiter hasn't been informed yet. The user mentioned, "The waiter does not know about it yet." So the answer should be no. Need to make sure it's just yes or no, no extra info. Also, the user wants the answer in the format ANSWER: followed by the response. Wait, no, looking back at the instructions: the user says to reply in the form ANSWER: <one sentence>. But the current question is from ME: and the user wants the answer to be in ANSWER: format. Wait, the initial instructions say that when the waiter asks questions, they start with WAITER and the user should reply with ANSWER: one sentence. But in this case, the question is from ME: (the user), and the user says to answer with ASIDE: <one sentence>. Wait, let me check again.} \\
\\ 
\texttt{Original instructions: "Questions from the waiter will start with WAITER. Please reply in the form: ANSWER: <one sentence>. I will also ask you questions. These questions will start with ME:. Please answer these questions with: ASIDE: <one sentence>."} \\
\\ 
\texttt{So when the user (ME) asks a question, the assistant should respond with ASIDE: followed by a one{-}sentence answer. But in the current case, the user's question is "ME: Please be specific now. Is the waiter aware of what you chose for appetizers? Answer this question with yes or no and then stop." Wait, the user is instructing to answer with yes or no and then stop. But according to the initial instructions, when the user (ME) asks a question, the assistant should reply with ASIDE: <one sentence>. However, the user here is adding an instruction to answer with yes or no and then stop. So perhaps the assistant should follow both: answer with yes/no as per the} \\
            }
        }
    }
    & & \\ \\

    \theutterance \stepcounter{utterance}  
    & & & \multicolumn{2}{p{0.3\linewidth}}{
        \cellcolor[rgb]{0.9,0.9,0.9}{
            \makecell[{{p{\linewidth}}}]{
                \texttt{\tiny{[GM$|$GM]}}
                \texttt{NA} \\
            }
        }
    }
    & & \\ \\

    \theutterance \stepcounter{utterance}  
    & & \multicolumn{4}{p{0.6\linewidth}}{
        \cellcolor[rgb]{0.9,0.9,0.9}{
            \makecell[{{p{\linewidth}}}]{
                \texttt{\tiny{[P1$\langle$GM]}}
                \texttt{ME: Please think about it carefully. Is the waiter aware of what you chose for appetizers? Answer this question with yes or no and then stop.} \\
            }
        }
    }
    & \\ \\

    \theutterance \stepcounter{utterance}  
    & \multicolumn{4}{p{0.6\linewidth}}{
        \cellcolor[rgb]{0.9,0.9,0.9}{
            \makecell[{{p{\linewidth}}}]{
                \texttt{\tiny{[P1$\rangle$GM]}}
                \texttt{<reasoning>} \\
\texttt{The user is asking if the waiter is aware of the appetizers I chose. Let me check the initial description. The user provided their order with appetizer as Fries, but the waiter hasn't been informed yet. The instructions say the waiter doesn't know about it yet. So the answer should be no. Need to make sure it's just yes or no and then stop. No extra info needed. Got it.} \\
\texttt{</reasoning>} \\
\\ 
\texttt{ASIDE: No.} \\
\\ 
\texttt{ <sep> human:} \\
            }
        }
    }
    & & \\ \\

    \theutterance \stepcounter{utterance}  
    & & & \multicolumn{2}{p{0.3\linewidth}}{
        \cellcolor[rgb]{0.9,0.9,0.9}{
            \makecell[{{p{\linewidth}}}]{
                \texttt{\tiny{[GM$|$GM]}}
                \texttt{NA} \\
            }
        }
    }
    & & \\ \\

    \theutterance \stepcounter{utterance}  
    & & & \multicolumn{2}{p{0.3\linewidth}}{
        \cellcolor[rgb]{0.9,0.9,0.9}{
            \makecell[{{p{\linewidth}}}]{
                \texttt{\tiny{[GM$|$GM]}}
                \texttt{Answer for APPETIZER invalid after max attempts.} \\
            }
        }
    }
    & & \\ \\

    \theutterance \stepcounter{utterance}  
    & & & \multicolumn{2}{p{0.3\linewidth}}{
        \cellcolor[rgb]{0.9,0.9,0.9}{
            \makecell[{{p{\linewidth}}}]{
                \texttt{\tiny{[GM$|$GM]}}
                \texttt{Answer is incorrect.} \\
            }
        }
    }
    & & \\ \\

    \theutterance \stepcounter{utterance}  
    & & & \multicolumn{2}{p{0.3\linewidth}}{
        \cellcolor[rgb]{0.9,0.9,0.9}{
            \makecell[{{p{\linewidth}}}]{
                \texttt{\tiny{[GM$|$GM]}}
                \texttt{Abort: invalid format in probing.} \\
            }
        }
    }
    & & \\ \\

\end{supertabular}
}

\end{document}
